\section{Results}

\subsection{Reference scores}

Table~\ref{tab:main} reports the seven steps at their primary agent profile. The full set
of seventeen step-and-profile combinations is released with the data.

Residual biorisk under baseline practice ranges from 3.33 (automated analysers, profile
P1) to 19.60 (open-tube culture work, profile P3). Two steps fall in the Elevated band or
above under baseline practice at their primary profile, and both involve open manipulation
of concentrated growth: MALDI-TOF target preparation at 17.73 (``High'') and open-tube
culture work at 19.60 (``Very high''). The recommended control sets move every step
downward, by a mean of 4.48 points across all seventeen combinations, but neither of the
two highest steps reaches the Low band even with the full recommended set applied: open
culture work falls to 13.07 and MALDI-TOF preparation to 10.27, both remaining Elevated.

\subsection{The inversion}
\label{sec:inversion}

The central result is visible in Fig.~\ref{fig:inversion}. Ordering the steps by residual
biorisk and plotting against evidence grade produces a strong negative relationship
between how much risk a step carries and how much is known about it.

The two highest-scoring steps are graded E3. For neither has aerosol release been measured
in a clinical laboratory. Their scores rest on incident investigation: sixteen cases of
probable laboratory-acquired meningococcal disease, of which eight were fatal, with the
isolate manipulated outside a biological safety cabinet in fifteen of the sixteen
laboratories~\cite{sejvar2005}; an estimated attack rate among microbiologists of 13 per
100{,}000 against 0.2 per 100{,}000 among adults in the general population, an
occupational rate roughly sixty-five times background~\cite{sejvar2005}; open-bench
subplating with a loop and open-bench aspiration from blood culture bottles without
cabinet, eye protection or mask among affected workers~\cite{mmwr2002}; and a fatal case
in which bench-top manipulation, open-bench plate spreading, and flaming of an
incompletely dried slide were among the breaches identified~\cite{mmwr2014}. For
MALDI-TOF target preparation, an investigation of \emph{Brucella} exposure risk events
across ten clinical laboratories found the identification technology in use during four
events, and those four accounted for 84\% of all personnel exposures identified across the
investigation~\cite{ackelsberg2020}.

The three steps graded E1 --- centrifugation, cap removal, and cartridge loading --- carry
baseline scores of 4.67, 6.33 and 10.27 respectively. The two pre-analytical steps are the
lowest-scoring in the reference at their primary profile. In both cases the measurement
was negative: air and surface sampling during centrifugation of spiked blood, and during
opening of spiked blood tubes, recovered no organisms, with an air sampling detection
limit of 4~CFU/m\textsuperscript{3}~\cite{dubuis2021}.

The highest-throughput step in the clinical laboratory, operation of automated chemistry
and haematology analysers, is graded E5. No study measuring aerosol release from such an
instrument was located. The only directly relevant statement in the literature is the
observation that automated laboratory equipment constitutes a source requiring
characterisation, studies having concentrated on manual procedures and
accidents~\cite{ravnholdt2025}.

\subsection{Where the release comes from}

Fig.~\ref{fig:items} decomposes the scores to the item level and shows that the two
highest-scoring steps reach their position by different routes.

Open-tube culture work scores at or near maximum on four of five release items: it applies
high localised energy through plate spreading, forceful streaking and flaming; it is open
by definition; it handles pure culture at maximum density, the most concentrated material
a clinical laboratory holds; and it is performed with the face close over the plate. Its
release profile is uniformly high.

MALDI-TOF target preparation, by contrast, scores low on energy. Smearing and spotting are
gentle manipulations. It reaches 17.73 through the combination of complete loss of
containment, enriched material, and close breathing-zone proximity, under the most severe
consequence profile. The hazard is not in the mass spectrometer; it is in the colony
handling upstream of it, and in the decision to identify an unrecognised isolate at all.

Centrifugation shows the opposite pattern: the highest energy score in the reference, with
containment, concentration and proximity all low. This is the structural reason it scores
as it does, and the reason the measured result was negative.

Panel (b) shows that baseline control reliability is low almost everywhere. Of the seven
steps, five score 1 on C4, reflecting that respiratory protection is not in routine use
for any of them; four score 1 on C2, reflecting controls that depend entirely on operator
memory. Only cartridge loading scores 4 on C1, because the closed cartridge is an
engineering control the laboratory receives without having to implement it.

\subsection{What the recommended controls are worth}

Fig.~\ref{fig:controls} shows the movement from baseline to recommended state. The
reductions are largest where baseline control is weakest and the consequence profile most
severe: open culture work falls 6.53 points and MALDI-TOF preparation 7.47 points, in both
cases driven principally by moving open manipulation into a certified Class~II cabinet.

This recommendation is supported by the only control-effectiveness observation in the
extracted evidence: across the \emph{Brucella} exposure events investigated, in those
events where a cabinet was in use for the manipulations concerned, no seroconversions were
identified among exposed personnel~\cite{ackelsberg2020}.

Cartridge loading falls 4.67 points, but the recommendation differs in character. The
closed cartridge is already a strong engineering control, and the measured evidence is
unusually favourable: neither sample preparation nor automated cartridge processing
produced any culturable bioaerosol, against 16 and 325~CFU/m\textsuperscript{3} measured
during preparation of smears for acid-fast microscopy~\cite{banada2010}. The recommended
set therefore protects the existing control --- correct reagent ratio, correct dispensing,
and correct contact time --- rather than adding new ones.

For centrifugation, the recommended set is sealed buckets, rotor inspection, and a written
breakage procedure. Respiratory protection is explicitly \emph{not} recommended, and C4 is
left at 1 in the recommended state, because it is not the control that addresses this
step's release mode.

\subsection{Sensitivity to the agent profile}

Fig.~\ref{fig:profiles} reports every step under every applicable profile. The spread is
wide: specimen reception ranges from 5.67 under P1 to 15.87 under P3, a difference of more
than two bands, from the same release and control scores. The profile assumption is
therefore the single largest determinant of where a step lands, which is the reason the
instrument reports three profiles rather than assuming one.

The practical reading is that the same physical step carries very different residual
biorisk according to what may be in the specimen, and that a laboratory which cannot
exclude an unsuspected low-dose inhalational agent is operating, for that step, in the
P3 column rather than the P1 column.

\subsection{The evidence map}

Across seven steps the grades distribute as: E1 in three, E2 in one, E3 in two, and E5 in
one. Three of seven steps therefore have no direct quantitative measurement of aerosol
release. Six of seven carry a conditionality flag, meaning that for all but one step the
evidence differs materially between the procedure performed as intended and the procedure
performed with a named failure or a technique deviation.

The single step without a conditionality flag is open-tube culture work. This is
substantively important: unlike the pre-analytical steps, it does not require a failure in
order to be hazardous. The documented infections occurred during the step performed as
intended, on an open bench.
